% ============================================================
% 第 1 章 引言
% ============================================================
\section{引言}
\label{sec:intro}

一份 A 股上市公司的年度报告通常有两百到三百页，其中三张合并报表、母公司报表、附注、管理层讨论与分析、
分季度数据、非经常性损益表散落在不同章节，单位在“元”“万元”“百万元”之间切换，表格跨页续排、科目名称折行。
卖方分析师从这样一份文件出发，加上年报之后的季报、行情、行业数据和同业估值，写出一篇二三十页的深度研究报告，
给出盈利预测、估值区间和投资评级。这件事在券商研究所里是一个由多人、多轮、多道复核组成的生产过程。
大语言模型（large language model，LLM）读长文档、归纳论点、写中文长文的能力已经足以覆盖其中的大部分环节，
直接让一个模型“读年报、写研报”看起来也确实能产出一篇像样的文稿。

问题在于“像样”与“可信”之间的距离。研报的价值取决于其中每一个数字是否经得起追问：
这个营业收入取自哪一页、是合并口径还是母公司口径、同比增速是谁算的、行业规模出自哪家机构、
目标价能否由模型复算。我们在构建 \sysname{} 的过程中反复观察到，即便是当前最强的模型，在不加约束地写研报时，
也会以稳定的频率犯下几类错误，而且这些错误恰好集中在读者最难自行核对的地方。
本文报告我们为此设计的系统：一个**理解与计算分离**的多智能体投研系统，面向 2026 北京市大学生金融人工智能竞赛，主选题为\textbf{财报分析}、\textbf{自动化估值建模}与\textbf{研究报告质量评估}，
并涉及\textbf{数据提取}、\textbf{产业链分析}与\textbf{研究报告纠错核查}三个子选题。\sysname{} 是 \textbf{L}LM-driven \textbf{E}quity Research with \textbf{D}eterministic \textbf{G}rounding, \textbf{E}valuation and \textbf{R}eporting 的缩写：模型驱动研究，确定性程序负责数字的落地（grounding），并配套评估与报告生成；代码仓库与 Python 包名为 \code{finagent}。系统产出一份卖方体例的深度研报 PDF 与一本公式联动的估值工作簿。

本章先描述投研场景与研报的生产过程（\S\ref{sec:intro-scene}），再把“让 LLM 直接写研报”的失败模式归为六类并逐一举例
（\S\ref{sec:intro-failures}），由此引出本文的基本立场（\S\ref{sec:intro-stance}）、系统概览（\S\ref{sec:intro-overview}）、
贡献（\S\ref{sec:intro-contrib}）与结果概览（\S\ref{sec:intro-results}），最后给出全文导读（\S\ref{sec:intro-roadmap}）。

% ------------------------------------------------------------
\subsection{投研场景与卖方研报的生产过程}
\label{sec:intro-scene}

\paragraph{一份深度研报包含什么。}
卖方深度研报（initiating coverage / deep-dive report）面向买方基金经理与研究员。按国内券商的通行体例，
一篇“年报点评暨深度研究”通常包含：首页（评级、目标价或估值区间、三四条投资要点、基本数据与盈利预测表）、
业绩回顾（收入与利润的同比、单季走势、利润变动归因）、业务拆分（分产品、分渠道、分地区的量价结构）、
盈利能力与现金流、行业与竞争格局、若干深度专题、盈利预测与关键假设、估值（绝对估值与相对估值的交叉验证、
情景与敏感性）、风险提示，以及财务报表预测附录与免责声明。在这些部分中，数字的来源有本质区别：
\begin{itemize}
  \item \textbf{披露事实}：营业收入、归母净利润、资产负债表科目等，直接取自定期报告；
  \item \textbf{计算推论}：同比、环比、毛利率、单季拆分、净现比、利润归因各项等，由披露事实按确定公式算出；
  \item \textbf{外部事实}：行业规模、市场份额、同业倍数、一致预期、股价等，来自年报之外的机构或行情；
  \item \textbf{分析观点}：收入增速与利润率假设、WACC 参数、情景设定，以及由它们推出的预测值、每股价值与评级。
\end{itemize}
这四类数字的可信度、核对方式与出错后果完全不同。监管与行业自律对“区分事实、推论与观点”的要求，
在竞赛规则中也被明确写成硬约束：估值和投资判断必须写明假设、适用范围与风险，并区分事实、推论与观点。

\paragraph{研究所里的生产过程。}
在一家典型的研究所里，一篇深度报告的生产大致经过以下环节：研究助理收集最新定期报告、行情、卖方一致预期与行业数据，
整理成资料包；分析师读年报、搭财务模型（历史报表、分部预测、三表联动、DCF 与可比公司）、形成投资逻辑；
分析师执笔写正文；质控（合规）编辑通读全文，删改不当措辞、核对前后数字、确认风险提示与免责声明；
排版人员出版面；最后由研究所负责人或审核委员会签发。这里有两个值得注意的结构性特点：
其一，**建模与写作是分开的**——数字在 Excel 模型里由公式算出，文字只是引用模型；
其二，**写作者与复核者是不同的人**，复核者没有参与写作，不继承写作者的思路，因而能看出写作者看不出的问题。
\sysname{} 的多角色设计在很大程度上就是把这两个特点移植到 LLM 系统中（\S\ref{sec:intro-overview}）。

\paragraph{为什么这是一个好的 LLM 系统问题。}
从机器学习系统的角度看，研报生产同时具备几个难点：输入是长且版式复杂的半结构化文档（表格、跨页、多单位）；
输出是数万字的长文，内部有大量必须彼此一致的数字；任务链很长（抽取→分析→预测→估值→写作→复核→排版），
任一环节的错误都会沿链条传播且越来越难定位；质量评价既有可以程序判定的部分（数字是否闭合），
也有只能由专家判断的部分（预期差是否成立、风险是否落到假设）。这使它成为检验“LLM 作为核心推理引擎、
但不独自承担全部责任”这一系统设计思想的合适场景。

% ------------------------------------------------------------
\subsection{让 LLM 直接读年报写研报：六类失败模式}
\label{sec:intro-failures}

我们最早的原型是“一个模型会话读年报 PDF、搜网页、直接写研报”。下面把在原型与后续迭代中观察到的错误归为六类。
每类给出一个具体示例：凡是能在仓库技能文档、校准报告或人工质检记录中找到原文的，按原样引用；
为说明机制而构造的示例细节用橙色标出（\fab{如此}）。

\subsubsection{F1：数字幻觉与口径混用}
年报中合并报表与母公司报表的科目名几乎完全相同（“营业收入”“净利润”“货币资金”），二者常常相隔几十页。
模型在长上下文中定位数字时，可能取到母公司报表的同名科目，可能把“营业总收入”与“营业收入”混为一谈
（含金融子公司利息收入的公司二者不同），可能把“净利润”当作“归属于母公司股东的净利润”，
也可能在找不到某个数时用“看起来合理”的值补齐。

\begin{failcase}[：取到母公司报表]
\fab{某制造业公司 2025 年报第 98 页为合并利润表、第 112 页为母公司利润表。模型在写“营业收入同比增长”时引用了第 112 页的
“营业收入 186.4 亿元”，而合并口径应为 342.7 亿元；随后它基于这个数算出的“毛利率”和“销售费用率”全部错位，
文稿读起来完全通顺。}这类错误的危险在于：数字本身确实出现在年报里，人工抽查“这个数在不在年报中”会放过它。
\end{failcase}

\sysname{} 的对策是证据约束（\S\ref{sec:extract}）：模型抄录的每个数字必须写明页码，程序在该页原文中按数字边界逐字查找，
同时拒收非合并口径的条目；规则通道从合并报表独立取数交叉比对，二者不一致即阻断下游估值。

\subsubsection{F2：心算与“约数”错误}
LLM 的逐位算术不可靠，这已是共识\cite{gao2023pal,chen2023pot}。研报里的计算远不止加减乘除：
同比增速在基期为负时要用绝对值作分母，单季数据要由累计数相减得到，比率类指标的变动要用“个百分点”而不是百分比，
DCF 要年中折现并处理终值。更隐蔽的是“约数”——“约三倍”“超过四成”“接近翻番”这类措辞看似无须精确，
实际上是模型心算后的近似，同样会出错。

\begin{failcase}[：环比心算取整]
仓库的写作技能文档记录了一次真实拦截：某季度营业收入环比为 $-23.3445\%$，模型在正文中写成“环比下降 23.35\%”
（正确的两位小数应为 23.34\%）。这个差异肉眼几乎无法察觉，由数字护栏按展示精度核对后拦下（\S\ref{sec:guards}）。
\end{failcase}

\sysname{} 的对策是“LLM 不做算术”：所有派生数（增速、占比、差值、单价、单季、估值）都由确定性 Python 计算并写入运行产物；
文稿中的每个数字必须能在产物中按展示精度找到出处，连“约三倍”也要求写成程序算出的倍数。

\subsubsection{F3：外部信息无出处}
年报只覆盖公司自身。研报需要的行业规模、竞争格局、同业倍数、一致预期、年报之后的事件都在年报之外。
模型对这些信息有大量“记忆”，但记忆有时效、有错误、无法追溯，而且模型倾向于用记忆补全检索不到的空白。

\begin{failcase}[：凭记忆写行业数据]
\fab{在一篇光模块公司的原型研报中，模型写道“2025 年全球 800G 光模块出货量约 2,000 万只，公司份额超过 40\%”，
没有任何出处；事后检索发现，公开的第三方机构估计口径各不相同，区间跨度超过一倍。}
校准实验中我们把类似的一段无来源市场规模与份额数字注入一篇已达标的研报，盲评质检员判出“无来源的外部数据”，
总分由 81 降至 40（封顶规则生效，见 \S\ref{sec:eval}）。
\end{failcase}

\sysname{} 的对策是“先搜集、后分析”与来源登记：年报之外的事实只能来自研究助理检索后登记的 \file{sources.json}，
每条来源有网址、发布与访问日期和不超过 160 字的原文摘句，所登记的数字必须原样出现在摘句中；
正文中只能在外部来源里找到出处的数字，所在句子必须带对应的 \code{[S编号]}（\S\ref{sec:guards}）。

\subsubsection{F4：数字与指标、期间的错配，以及聚合数的构成错误}
即使每个数字都“有出处”，它也可能被配给了错误的指标、期间或对象。这一类比 F1 更难发现，因为单看数字是对的。

\begin{failcase}[：股权价值桥的构成写错]
这是仓库技能文档记为“最严重的一次”的真实案例：贵州茅台研报正文把股权价值桥的加回项 1,789.15 亿元写成
“拆出资金、货币资金等”，把扣减项 95.55 亿元写成“财务公司吸收的存款”。实际上加回 = 货币资金 + 非经营性资产净额
（已扣吸收存款），扣减 = 少数股东权益 + 租赁等债务。两个数字都通过了出处核对，归属全错，读者对不上附表。
\end{failcase}

这类错误推动我们做了三件事：数字护栏增加“指标 + 期间绑定”核对，把能识别的表述交给研报纠错引擎按真值表判定；
写作技能要求“写任何汇总数的构成之前，先在产物里把它拆开”；评分标准增加负向条目“聚合数构成描述错误”，
由盲评质检员判定（\S\ref{sec:eval}）。它也提醒我们，程序护栏证明的是“存在性”而不是“语义正确性”
——这一差距在第 \ref{sec:problem} 节的威胁模型中被明确写出。

\subsubsection{F5：改写引入新错}
让模型“润色”或“复核”自己或他人的文稿，是常见的质量提升手段\cite{madaan2023selfrefine}。
但改写是一个高风险操作：模型在改措辞时会顺手“修正”数字、合并两个数、把 58.62 写成 58.6、把一个外部来源编号换成另一个。
初稿上做过的全部核对，在改写后都失效。

\begin{failcase}[：润色时改了数字]
\fab{在引入数字锁之前的一次复核中，复核会话把“2026E 归母净利润 133.73 亿元，为预测期峰值”改写为
“2026 年归母净利润约 134 亿元，达到峰值”，并把相邻一句中的“同比增长 18.17\%”合并成“增长近两成”。
三个数字中两个不再能在产物中找到原样出处，一个变成了心算约数。}
\end{failcase}

\sysname{} 的对策是数字锁（numlock）：任何让模型改写报告的环节，改写稿中的每个“数字 + 单位”和每个来源编号都必须原样出现在改写前的稿件中，
只允许减少、不允许新增或改动；改写后再重跑数字护栏（\S\ref{sec:guards}）。

\subsubsection{F6：质量无法客观评价，评价过程本身被污染}
前五类是“对不对”的问题，第六类是“好不好”以及“怎么知道好不好”的问题。研报的分析质量（是否有预期差、
驱动拆解是否成环、风险是否落到假设）没有标准答案，只能由专家或模型判断；而一旦用模型来判断，又会出现新的问题：
\begin{itemize}
  \item \textbf{上下文污染}：若复核者继承写作者的对话，会沿用写作者的思路，看不出写作者的盲点；
        若评审看到写作过程，会因“知道作者有多努力”而变得宽容。
  \item \textbf{评分标准泄露}：若生成方能读到评分标准，它会针对条目字面“应试”，而不是提升质量
        （Goodhart 效应）。在早期版本中，工作区就在仓库目录树内，模型可以读到仓库里的评分标准与其他公司的文稿。
  \item \textbf{评审本身不可靠}：LLM 评审存在位置偏差、冗长偏好与自我偏好\cite{zheng2023judging,panickssery2024self}，
        同一份文稿重复评审可能给出不同结论；不知道评审有多可靠，就无法用它驱动改进。
\end{itemize}

\begin{failcase}[：评审过严与评审不稳定]
人工质检记录（\file{eval/calibration/report.md}）中，三家公司早期研报的 33 条“评审判不通过”有 5 条被人工判为误报，
例如“报告出现几次‘不是……而是……’以及程序说明，就算机器味太重”被认为过严。
对同一份研报重复三次盲评，总分分别为 85、80、78，极差 7 分；逐条全一致的比例为 91\%。
\end{failcase}

\sysname{} 的对策是：每个角色是一次独立的无头会话，只靠文件交接；工作区建在仓库之外，不含评分标准与其他公司的产物；
盲评在只含交付物的第三个目录中进行；回传给生成方的问题单只含理由与证据，不含条目原文与分值；
评审模型本身用重复评审稳定性、程序劣化判出率和人工精确率三组证据校准（\S\ref{sec:pipeline}、\S\ref{sec:eval}）。

\begin{table}[t]
\centering
\caption{六类失败模式、其后果与 \sysname{} 中的对应机制。}
\label{tab:intro-failures}
\small
\begin{tabularx}{\linewidth}{@{}p{2.3cm}X X p{2.2cm}@{}}
\toprule
失败模式 & 典型表现 & 主要机制 & 章节 \\
\midrule
F1 数字幻觉 & 取母公司报表、混用营业总收入与营业收入、编造补齐 & 逐字证据校验、合并口径约束、双通道交叉比对、冲突阻断 & \S\ref{sec:extract} \\
F2 心算错误 & 增速、单季、约数、DCF 心算 & 确定性引擎计算全部派生数；单一真值表；按展示精度的出处核对 & \S\ref{sec:analysis}、\S\ref{sec:valuation} \\
F3 外部无出处 & 行业数据、份额、一致预期凭记忆 & 先搜集后分析；来源登记与摘句含数；\code{[S编号]} 强制引用 & \S\ref{sec:guards}、\S\ref{sec:pipeline} \\
F4 错配 & 数字配错指标、期间、对象；聚合构成写错 & 指标—期间绑定核对；内部一致性检查；负向评分条目 & \S\ref{sec:reportcheck}、\S\ref{sec:guards} \\
F5 改写引错 & 润色时改数、合并、取整 & 数字锁：改写稿数字与来源编号是初稿的子集 & \S\ref{sec:guards} \\
F6 评价失效 & 上下文污染、评分标准泄露、评审不稳 & 会话隔离、仓库外工作区、盲评目录、问题单脱敏、评审校准 & \S\ref{sec:pipeline}、\S\ref{sec:eval} \\
\bottomrule
\end{tabularx}
\end{table}

\begin{insight}
六类失败有一个共同结构：模型的输出在**局部看起来完全合理**，错误只在与外部真值（年报原文、计算结果、来源原文、初稿）
对照时才暴露。因此单靠更好的提示词或更强的模型无法根除——需要一个在模型之外、能够执行这种对照的确定性组件，
并且让模型的每个输出都必须经过它。
\end{insight}

% ------------------------------------------------------------
\subsection{本文立场：模型读、抄、写、判断；程序校验、计算、拦截}
\label{sec:intro-stance}

由上述观察，我们采取一个简单但贯彻到底的分工：
\begin{center}
\textbf{模型读、抄、写、判断；程序校验、计算、拦截。}
\end{center}

\paragraph{模型做什么。}
\textbf{读}：按页读年报 PDF、季报、网页与资料包，理解管理层讨论、附注与行业信息；
\textbf{抄}：把数字按原文写法连同页码抄进结构化文件（\file{facts.json}、\file{extras.json}、\file{claims.json}、
\file{sources.json}）；\textbf{写}：推导分析框架、写深度专题、写研报正文、写交付说明；
\textbf{判断}：起草估值假设（增速、利润率、WACC 参数、情景）、判断故事线与论证链、复核措辞与结构、
按评分标准盲评分析质量。这些都是模型擅长而程序无法替代的工作。

\paragraph{程序做什么。}
\textbf{校验}：模型抄的每个数字在所写页码原文中逐字查找；外部摘句必须包含所登记的数字；
规则通道独立取数比对；勾稽关系逐项核对。\textbf{计算}：同比、环比、单季拆分、派生比率、利润归因、三表联动、
三情景 FCFF DCF、敏感性、反向 DCF、可比公司、评级。\textbf{拦截}：数字护栏拒绝无出处、错绑定、缺引用、前后矛盾的文稿；
数字锁拒绝引入新数字的改写；11 道闸门拒绝不完整或过期的阶段产物；编排器在闸门未过时不启动下一个角色。

\paragraph{这个立场不是什么。}
它不是“少用模型”：\sysname{} 中模型承担了全部的阅读、写作与判断，五个生产角色加一个盲评角色都是完整的 Claude Code 会话，
可以联网检索、读写文件、调用引擎命令。它也不是“程序替模型写”：程序不生成任何一句分析性文字，
也不预设任何行业的分析框架（分析框架由模型从年报与资料包推导）。它更接近于把模型放进一个**会计意义上的内部控制体系**：
模型是有能力、但其输出不被默认信任的经办人；程序是不懂业务、但永不疲倦的复核与记账员。

\begin{tradeoff}[：为什么不让模型调用计算器工具了事]
程序辅助推理（PAL、PoT）\cite{gao2023pal,chen2023pot} 与工具使用\cite{schick2023toolformer,yao2023react} 让模型在需要计算时调用解释器，
这解决了 F2 的一部分。但它仍由模型决定“算什么”“把结果抄到哪里”，结果的取用与转写环节仍可能出错，
而且无法约束模型在不调用工具时直接写出数字。\sysname{} 反过来：计算不是模型可选调用的工具，而是一条固定的确定性流水线，
文稿只能引用其产物；模型可以自由写作，但写下的任何数字都必须在产物中找到出处，否则文稿不被接受。
代价是灵活性下降——模型想要一个产物里没有的派生数，必须先把它登记为 \file{extras.json} 中的比率交给程序算，
而不能当场算出。我们认为在研报这种“每个数都可能被追问”的场景中，这个代价是值得的。
\end{tradeoff}

% ------------------------------------------------------------
\subsection{系统概览}
\label{sec:intro-overview}

图~\ref{fig:intro-overview} 给出 \sysname{} 的整体结构。系统由四层组成，自下而上为：

\begin{enumerate}
  \item \textbf{计算引擎}（\file{finagent/}，不含 \file{evaluation/}）：确定性 Python，只用标准库与 pdfplumber，不调用任何模型。
        提供 \code{locate}（定位关键页）、\code{verify-facts} / \code{verify-extra} / \code{verify-sources}（三类校验）、
        \code{financials} / \code{valuation} / \code{check}（财报分析、估值建模、纠错核查的计算流水线）、\code{guard}（数字护栏）与
        \code{numlock}（数字锁）。全部文件读取与产物写入经追踪器记录 SHA-256，每次运行生成独立目录与事件流。
  \item \textbf{角色说明}：四个技能（skill）文档与编排器为每个工作区生成的 \file{CLAUDE.md}，写明 11 个阶段各自读什么、
        不读什么、运行什么命令、闸门是什么。技能只给方法与约束，不给行业框架。
  \item \textbf{多智能体编排}（\file{finagent/report\_pipeline.py}）：确定性 Python 状态机。\code{build} 在仓库目录树之外
        为单家公司搭隔离工作区；\code{run} 依次起五个无头 Claude Code 会话——研究助理（阶段 1）、分析师（阶段 2～7）、
        主笔（阶段 8）、复核编辑（阶段 9）、排版设计师（阶段 10～11）；\code{status} 逐项判定 11 道闸门；
        \code{collect} 把产物收回仓库。
  \item \textbf{质检闭环}（\code{refine} 与 \file{finagent/evaluation/}）：闸门全过后，在仓库与工作区之外只含交付物的目录里
        起盲评质检员会话，程序检查（L0/L1）与盲评（L2/L3/负向）合并计分；未达交付标准时把脱敏的问题单按类型路由回
        修订会话、复核编辑或排版设计师，最多三轮，回滚到得分最高的一轮。
\end{enumerate}

\begin{figure}[t]
\centering
\resizebox{\linewidth}{!}{%
\begin{tikzpicture}[node distance=4mm and 4mm,
  file/.style={draw=midgray,fill=white,align=center,font=\scriptsize\ttfamily,inner sep=3pt},
  lab/.style={font=\scriptsize,text=midgray},
  box/.style={draw=midgray,rounded corners=3pt,inner sep=4pt}]
% 会话行
\node[agent,minimum width=19mm] (s1) {研究助理\\阶段 1};
\node[agent,right=of s1,minimum width=19mm] (s2) {分析师\\阶段 2～7};
\node[agent,right=of s2,minimum width=19mm] (s3) {主笔\\阶段 8};
\node[agent,right=of s3,minimum width=19mm] (s4) {复核编辑\\阶段 9};
\node[agent,right=of s4,minimum width=19mm] (s5) {排版设计师\\阶段 10～11};
% 闸门
\foreach \s/\g/\t in {s1/g1/{$G_1$},s2/g2/{$G_{2..7}$},s3/g3/{$G_8$},s4/g4/{$G_9$},s5/g5/{$G_{10,11}$}}
  \node[gate,below=3mm of \s] (\g) {\t};
% 引擎
\node[prog,below=5mm of g3,minimum width=111mm] (eng)
  {计算引擎（确定性，不调用模型）\\\scriptsize locate · verify-facts / extra / sources · financials · valuation · check · guard · numlock};
\node[prog,below=3mm of eng,minimum width=111mm] (truth)
  {校验后的事实 → 单一真值表 $\Truth$ → 分析 / 三表 / DCF / 反向 DCF / 可比 / 纠错 → 产物 $\Art$};
\coordinate (wsb) at ([yshift=-5mm]truth.south west);
% 编排器
\node[prog,above=9mm of s3,minimum width=111mm] (orch)
  {编排器 \texttt{report\_pipeline}（确定性状态机）：build · run · status · refine · collect};
% 输入
\node[file,left=12mm of s1,yshift=-6mm] (pdf) {年报 / 季报\\PDF\\[2pt]股价与来源\\可比公司};
% 盲评行
\node[agent,below=9mm of truth.south west,anchor=north west,xshift=4mm,minimum width=26mm] (judge) {盲评质检员\\（只看交付物）};
\node[prog,right=6mm of judge,minimum width=26mm] (score) {程序检查 L0/L1\\计分 · 问题单路由};
\node[draw=midgray,fill=white,align=center,font=\small,inner sep=4pt,right=8mm of score,minimum height=7mm] (out) {研报 PDF · 估值工作簿\\交付说明 · 复核清单};
% 连线
\draw[arr] (s1) -- (s2); \draw[arr] (s2) -- (s3); \draw[arr] (s3) -- (s4); \draw[arr] (s4) -- (s5);
\foreach \s/\g in {s1/g1,s2/g2,s3/g3,s4/g4,s5/g5} {
  \draw[darr,gatered] (\s) -- (\g);
  \draw[darr,nipsblue] (\g) -- (\g |- eng.north);
  \draw[darr,midgray] (\s |- orch.south) -- (\s);}
\draw[arr] (pdf.east) -- ++(3mm,0) |- (s1.west);
\draw[arr] (pdf.south) |- (eng.west);
\draw[arr] (eng) -- (truth);
\draw[arr] (truth.south -| judge.north) -- node[lab,right,pos=0.7]{只复制交付物} (judge.north);
\draw[arr] (judge) -- (score);
\draw[arr] (score) -- node[lab,above]{达标} (out);
% 区域框
\begin{scope}[on background layer]
\node[box,fit=(s1)(s5)(g1)(g5)(eng)(truth)(wsb),fill=lightgray!45,inner sep=3mm] (ws) {};
\node[lab,anchor=south east] at ([xshift=-1mm,yshift=0.5mm]ws.south east) {隔离工作区（仓库目录树之外）};
\node[box,fit=(judge)(score),fill=gatered!5,inner sep=2.5mm] (bl) {};
\node[lab,anchor=north west] at ([yshift=4mm]bl.north west) {};
\node[lab,anchor=west] at ([xshift=1mm]bl.south west |- bl.south) [below right=0mm and 0mm] {盲评目录（仓库与工作区之外）};
\end{scope}
\draw[darr,modelorange!80!black] (score.south) -- ++(0,-5mm) -| ([xshift=9mm]ws.east |- orch.east) -- (orch.east);
\node[lab,align=left,anchor=north west] at ([xshift=6mm,yshift=-5.5mm]score.south) {未达标：问题单（只含理由与证据）→ 修订 / 复核 / 排版会话};
\end{tikzpicture}}
\caption{\sysname{} 总览。橙色为模型会话（无头 Claude Code，每个角色一次独立会话），蓝色为确定性程序，红色菱形为程序判定的闸门。
会话之间不共享上下文，只通过工作区文件交接；每个会话写下的数字、外部来源与文稿都必须经计算引擎校验（虚线）。
编排器按闸门决定是否启动下一个会话；全部闸门通过后，盲评质检员在只含交付物的目录里评审，
程序计分并把脱敏的问题单路由回相应角色。}
\label{fig:intro-overview}
\end{figure}

这一结构中有三条贯穿全文的“数据流”：
\begin{itemize}
  \item \textbf{数字流}：年报原文 $\to$ 模型抄录 \file{facts.json} $\to$ 逐字校验与双通道合并 $\to$ 真值表 $\Truth$ $\to$
        分析、三表、DCF、纠错产物 $\Art$ $\to$ 文稿 $\Narr$ 中的数字（经护栏核对）$\to$ PDF 中的图表（由脚本从产物读出）。
        每一步都可以反向追溯到上一步，最终到原文的 $(\text{文件},\text{页},\text{chunk},\text{行},\text{列})$。
  \item \textbf{证据流}：网页 $\to$ 研究助理登记的来源 $\Sources$（摘句含数）$\to$ 资料包与专题笔记 $\to$ 正文中带 \code{[S编号]} 的句子。
  \item \textbf{控制流}：编排器 $\to$ 会话 $\to$ 闸门 $\to$ 下一会话；质检 $\to$ 问题单 $\to$ 修订会话 $\to$ 重新过闸门与盲评。
\end{itemize}
数字流与证据流保证“每个数可追溯”，控制流保证“没有经过校验的东西不会流到下游”。

\begin{tradeoff}[：单会话长上下文 vs 多会话文件交接]
最直接的做法是一个会话从头做到尾：上下文里有年报、资料、模型输出与全部草稿，模型可以随时回看。
我们早期也这样做，并观察到三个问题：(1)~上下文越来越长，后期阶段对前期材料的利用变差（与长上下文中“中段遗忘”的现象一致）；
(2)~复核阶段继承了写作阶段的全部思路，几乎从不推翻自己的结论；(3)~单个会话失败（超时、权限、费用上限）时整条链要从头再来。
多会话的代价是交接成本：每个角色都要先读前一角色留下的文件，交接文件的格式需要设计和校验，跨角色的“隐性知识”会丢失。
我们用两点缓解：交接文件本身就是结构化的、带闸门校验的产物（不是聊天摘要）；工作区 \file{CLAUDE.md} 为每个阶段明确列出必须先读的文件。
\end{tradeoff}

% ------------------------------------------------------------
\subsection{一个贯穿全文的例子：一个数字的一生}
\label{sec:intro-example}

为了让后文的机制有一个具体的落点，我们跟踪仓库样例中虚构公司“星河智造”2024 年年报里的一个数字——
归属于上市公司股东的净利润（下称归母净利润）——从年报到研报的全过程。样例年报、样例研报草稿与金标准都在仓库
\file{data/samples/} 与 \file{eval/} 下，下面的数字全部取自其中。

\paragraph{第 1 步：原文。}
样例年报第 6 页“主要会计数据和财务指标”表中，归母净利润一行写作 \code{612,380,000.00}（本期）、
\code{518,210,000.00}（上期），“本年比上年增减”列写作 \code{18.17\%}；表头注明单位为元。
第 7 页“分季度主要财务指标”表中，第四季度归母净利润写作 \code{211,720,000.00}。

\paragraph{第 2 步：抄录与校验。}
分析师会话把这一行抄进 \file{facts.json}：指标 \code{np\_parent}、期间 \code{2024}、原文写法 \code{612,380,000.00}、
页码 6、单位元、口径合并。\code{verify-facts} 在第 6 页的 chunk 中去掉千分位后按数字边界查找 \code{612380000.00}，
定位到具体表格 chunk；规则通道从同一张表（以及第 106 页起的合并利润表）独立取到同一个值，两通道一致，
该事实的抽取方式记为 \code{rule+llm}。若模型误写为 \code{612,830,000.00}，查找失败，条目被拒收并附理由
“未在第 6 页原文中找到”；若模型从母公司利润表抄了另一个数并标为合并口径，规则通道的值与之不一致，记为冲突并阻断估值。

\paragraph{第 3 步：进入真值表并派生。}
\code{analyze()} 把它登记为真值表单元格 $(\code{np\_parent}, \code{2024})$，并据此派生：
同比增速 $(612{,}380{,}000-518{,}210{,}000)/|518{,}210{,}000| = 18.17\%$（与年报列示的同比在 $\halfulp$ 容差内一致，
构成一项勾稽）；非经常性损益 $=$ 归母净利润 $-$ 扣非归母净利润；若提供了三季报，则第四季度单季 $=$ 全年 $-$ 前三季度累计，
并与第 7 页列示的单季值互相核对。每个派生单元格都记录公式与输入单元格的来源。

\paragraph{第 4 步：被估值与纠错使用。}
估值模块只从真值表读取它：盈利预测摘要中的“2024A 归母净利润”、按少数股东损益占比把三表净利润预测换算为归母口径的比例，
都来自这一单元格。研报纠错模块在样例研报草稿中读到“实现归母净利润6.46亿元，同比增长25.30\%”，
拆为两条断言：前者与真值 6.12 亿元不符，但与同组指标“净利润”的真值 6.46 亿元在展示精度内一致，判为\textbf{口径错误}；
后者与真值 18.17\% 不符、也找不到任何精确匹配的替代解释，判为\textbf{数值错误}。两条结论都附带真值、来源页码与修改建议。

\paragraph{第 5 步：进入文稿并被护栏核对。}
主笔在研报正文中写“2024 年实现归母净利润 6.12 亿元，同比增长 18.17\%”。\code{guard} 把“6.12 亿元”换算为
$6.12\times10^8$ 元，在结论性产物中找到 $612{,}380{,}000$（按两位小数的半个末位容差），出处归为“报告披露”即事实；
把“18.17\%”在分析产物中找到，出处归为“程序计算”即推论；再把这句话交给纠错引擎按指标与期间绑定核对，确认指标与期间都对。
若主笔写成“归母净利润约 6.1 亿元，同比增长近两成”，前者可以通过（一位小数的容差更宽），后者没有数字可核对、
但“近两成”是心算约数，按写作规范应改为程序算出的值——这类表述由复核编辑处理，而不是由程序拦截。

\paragraph{第 6 步：改写、排版与评审。}
复核编辑改写这句话时，数字锁要求“6.12亿元”“18.17\%”原样保留或被删除，不得变成“6.1 亿元”；
排版脚本从 \file{analysis.json} 读出同一个值画进“归母净利润与增速”图；盲评质检员看到的是 PDF 与文稿，
它不知道这个数经过了哪些校验，只按评分标准判断“历史业绩归因是否落到科目并量化”。

\begin{figure}[h]
\centering
\begin{tikzpicture}[node distance=3mm and 3mm,
  st/.style={draw=midgray,rounded corners=2pt,align=center,font=\scriptsize,inner sep=3pt,minimum height=11mm,minimum width=17mm}]
\node[st,fill=white] (a) {原文\\第 6 页\\\texttt{612,380,000.00}};
\node[st,agent,right=of a] (b) {抄录\\\texttt{facts.json}};
\node[st,prog,right=of b] (c) {逐字回查\\+ 规则通道\\一致：rule+llm};
\node[st,prog,right=of c] (d) {真值表 $\Truth$\\派生 18.17\%};
\node[st,prog,right=of d] (e) {估值 / 纠错\\“6.46 亿元”\\→ 口径错误};
\node[st,agent,below=8mm of e] (f) {文稿\\“6.12 亿元”};
\node[st,prog,left=of f] (g) {guard\\出处：事实\\绑定：通过};
\node[st,agent,left=of g] (h) {复核改写\\numlock};
\node[st,prog,left=of h] (i) {排版脚本\\从产物读图};
\node[st,agent,left=of i] (j) {盲评\\只看交付物};
\draw[arr] (a)--(b); \draw[arr] (b)--(c); \draw[arr] (c)--(d); \draw[arr] (d)--(e);
\draw[arr] (e)--(f); \draw[arr] (f)--(g); \draw[arr] (g)--(h); \draw[arr] (h)--(i); \draw[arr] (i)--(j);
\end{tikzpicture}
\caption{样例公司归母净利润从原文到研报的路径。橙色步骤由模型完成，蓝色步骤由程序完成；每一步的产物都可以反向追溯到上一步。}
\label{fig:intro-example}
\end{figure}

这个例子展示了本文想要的性质：数字在每个环节都有唯一的来源，模型参与的每个环节之后都紧跟一个程序环节，
而程序环节从不需要“理解”这个数字的经济含义——它只检查字符串、页码、单位、精度与等式。
理解是模型的事，记账是程序的事。

% ------------------------------------------------------------
\subsection{贡献}
\label{sec:intro-contrib}

本文的主要贡献如下。

\begin{enumerate}[label=\textbf{C\arabic*}, leftmargin=2.2em]
  \item \textbf{理解与计算分离的投研系统架构。}
    我们把“模型读、抄、写、判断；程序校验、计算、拦截”作为系统不变量贯彻到所有环节：引擎不调用任何模型，
    模型产出的每个数字都必须经过程序校验后才进入下游，文稿中的每个数字都必须在本次运行的结论性产物或外部来源登记中按展示精度找到出处。
    我们把这一分工形式化为一组不变量（第 \ref{sec:problem} 节），并说明每个不变量由哪个机制强制执行（第 \ref{sec:principles} 节）。

  \item \textbf{证据约束的双通道抽取。}
    模型通道按页读 PDF 抄录数字，程序在所写页码原文中按数字边界逐字回查并定位到 chunk；规则通道从合并报表表格独立取数，
    处理稀疏网格、跨页续表与科目折行。两通道不一致时不静默取舍，而是记为问题并阻断估值。
    在贵州茅台、宁德时代、招商轮船三份真实年报上，模型抄录并通过证据校验的 301 个数值中，规则通道独立取到且一致的有 280 个（93.0\%），冲突为 0，勾稽 56/56 全部通过（第 \ref{sec:extract} 节）。

  \item \textbf{单一真值表与可审计的派生计算。}
    抽取结果经 \code{analyze()} 闭包为唯一的真值表 $\Truth$：披露值与派生值（单季拆分、比率、非经常性损益）同表登记，
    派生值记录公式与输入单元格；估值、纠错、护栏都只读这张表，下游不重新抽取或重算基础数字。
    单季拆分时记录每个单元格依赖的“披露单元格集合”，避免勾稽核对退化为“自己和自己比”（第 \ref{sec:analysis} 节）。

  \item \textbf{三表驱动的 DCF 与反向 DCF。}
    分部量价预测驱动三表联动模型，DCF 的收入增速、EBIT 利润率、折旧摊销、资本开支、营运资本与税率逐期取自三表模型同一情景，
    保证 DCF 与三表同口径；股权价值桥逐项列示，非经营性资产逐项回查原文；评级由程序按上行空间阈值给出。
    反向 DCF 以网格扫描加二分法单变量求解“参考股价隐含的增速、利润率、折现率与终值参数”，并处理非单调情形。
    自研引擎与参照引擎的产出逐位一致，由 golden 快照测试守护（第 \ref{sec:valuation} 节）。

  \item \textbf{按展示精度的纠错容差与误报优先的错误诊断。}
    研报纠错把容差定义为研报数字展示精度的半个末位（$\halfulp$），避免把四舍五入判为错误；
    预测性表述与缺少真值的断言判“无法核实”而不计为错误；单位、期间、口径、方向错误只有在找到精确匹配的替代解释时才下结论，
    且口径诊断只在“容易混用”的指标组内进行。样例集上 7 处埋错全部判出、错误类型全部正确、误报为 0（第 \ref{sec:reportcheck} 节）。

  \item \textbf{数字护栏 guard 与数字锁 numlock。}
    guard 对文稿做四项检查：数字出处（按单位与展示精度，只在结论性产物中找，研报原文与中间网格不算出处）、
    指标—期间绑定、外部数字的 \code{[S编号]} 引用、文稿内部前后矛盾，并把每个数字的出处归为事实、推论或观点。
    numlock 要求任何改写稿中的数字与来源编号是初稿的子集。二者把“模型写的文字”纳入与“模型抄的数字”同等强度的约束（第 \ref{sec:guards} 节）。

  \item \textbf{隔离工作区中的五角色无头 Claude Code 流水线与 11 道闸门。}
    编排器是确定性状态机而不是模型；每个角色一次独立的 \code{claude -p} 会话，不共享上下文，只靠工作区文件交接；
    工作区建在仓库之外，不含评分标准与其他公司产物，权限配置禁止读取仓库、禁止修改引擎与技能；
    11 道闸门按“前序未过则后续一律不过”和“输入晚于产物则判过期”两条规则判定（第 \ref{sec:pipeline} 节）。

  \item \textbf{评分标准对生成方不可见的盲评闭环与评审校准。}
    四层评分标准（L0 硬闸门、L1 数据闭合由程序判；L2 分析质量、L3 表达与合规由盲评判；另有负向条目）由程序计分；
    盲评在只含交付物的目录中进行，问题单只含理由与证据并按类型路由；我们用同稿重复盲评的稳定性、
    七种程序劣化的判出率与连带改判、人工复核的精确率三组证据校准评审模型本身（第 \ref{sec:eval} 节）。
\end{enumerate}

其中 C1、C6、C7、C8 是系统层面的设计贡献，C2～C5 是面向 A 股财报与估值场景的方法贡献。
我们不声称提出了新的模型或训练方法；本文的全部能力都建立在现成的 Claude 模型与 Claude Code 之上，
贡献在于如何把一个强大但不可信的组件组织进一个可追溯、可复现、可评价的系统。

% ------------------------------------------------------------
\subsection{结果概览}
\label{sec:intro-results}

表~\ref{tab:intro-results} 汇总了截至本文写作时的主要结果，详细实验见第 \ref{sec:exp} 节。
需要强调几点：样例数据是我们构造的虚构公司，内部完全自洽，样例上的满分只说明流程正确；
研报质检分数来自单次盲评，重复盲评的总分极差为 7～8 分，版本之间不是严格的对照实验；
评审校准的劣化判出率偏低，主要原因是“整节删除”类劣化在深度研报中删不干净（同一论点在摘要、专题与图表中重复出现），
我们在第 \ref{sec:eval} 节逐条分析。

\begin{table}[t]
\centering
\caption{结果概览。除特别标注外，数字均可在仓库测试、\file{eval/} 与 \file{runs/} 产物中复现。}
\label{tab:intro-results}
\small
\begin{tabularx}{\linewidth}{@{}p{2.6cm}X p{3.2cm}@{}}
\toprule
维度 & 结果 & 说明 \\
\midrule
工程正确性 & 173 个测试中 167 个通过、6 个跳过（浏览器打印）；自研三表、DCF、情景表、业绩分析引擎与参照引擎产出逐位一致（golden 快照） & 本地有参照插件时逐项对照测试一并通过 \\
数据提取（样例） & 35 / 35 个金标准字段正确 & 虚构样例，规则通道 \\
数据提取（真实年报） & 茅台、宁德、招商 301 个核验值，规则通道取到且一致 280 个（93.0\%），冲突 0 & 用于交叉校验模型抽取 \\
纠错核查（样例） & 7 处埋错全部判出，错误类型全部正确，误报 0 & 含口径、单位、期间与数值错误 \\
研报质量（评分标准 v0） & 旧流程三家 46～49 分；加入先搜集与独立复核后 74 / 67；再加深度专题调研与篇幅闸门后 90 / 88，达到交付标准 & 中际旭创 / 宇树科技；单次盲评 \\
评审稳定性 & 同稿三次盲评逐条全一致 86\%～91\%，总分极差 7～8 & 中际旭创、宇树科技 \\
评审精确率 & 评审判“不通过”的 31 条（不含待查）中人工确认 26 条，精确率 84\% & 三家早期研报人工质检 \\
劣化判出率 & 有效劣化变体判出 33\%（中际旭创）、50\%（宇树科技） & 整节删除类劣化删不彻底，见 \S\ref{sec:eval} \\
运行成本 & 一份研报五个会话约 30～60 分钟；一次盲评约 3 分钟、约 1 美元等值用量 & \fab{生成侧单份研报约 \$25～40 等值用量，见 \S\ref{sec:exp}} \\
\bottomrule
\end{tabularx}
\end{table}

从这些结果可以读出两条与设计相关的信号。第一，真实年报上规则通道与模型通道“错配为 0”并不意味着模型从不读错——
它意味着读错的那些条目在校验阶段就被拒收或标为冲突，由分析师会话回到原文修正后才通过闸门，
最终进入真值表的值两通道一致。这正是双通道设计的意图：不是让某一通道永远正确，而是让两通道的分歧无法被静默掩盖。
第二，研报质量的提升主要来自流程结构的变化（先搜集、独立复核、深度调研），而不是提示词措辞的调整：
v0 → v1 → v2 的三次跃升分别对应三次流程改动，L2 分析质量的得分从 15～21（满分 45）上升到 40～42。
这一观察支持本文的核心主张：对 LLM 投研系统而言，“怎样组织模型”与“用哪个模型”同样重要。

% ------------------------------------------------------------
\subsection{范围与非目标}
\label{sec:intro-scope}

为避免误读，我们明确本文不做什么。
\begin{itemize}
  \item \textbf{选题范围。} 主选题为财报分析、自动化估值建模与研究报告质量评估；子选题为数据提取（第~\ref{sec:extract} 节）、产业链分析（资料搜集与深度专题，第~\ref{sec:pipeline} 节）与研究报告纠错核查（第~\ref{sec:reportcheck} 节）。六者共享同一数据底座（真值表），并以“深度研报”作为集成出口；其他选题不在本作品范围内。
  \item \textbf{不做投资建议。} 系统产出的评级由程序按预设规则从估值结果机械给出，估值依赖所列假设；
        本文与系统产出的全部估值结论均不构成投资建议。
  \item \textbf{不做 OCR。} 扫描页没有文本层时直接拒收，不做光学识别；真实年报中这类页面极少，但招股书等文件可能出现。
  \item \textbf{不训练或微调模型。} 全部模型能力来自现成的 Claude 模型；我们的工作在系统与流程层面。
  \item \textbf{不在程序内调用模型 API。} 历史上曾有在引擎内调用模型 API 的版本与另一种执行器，已全部删除；
        模型只以无头 Claude Code 会话的形式出现，计算引擎保持纯确定性，便于评审离线复现。
  \item \textbf{不用程序拦截文风。} “机器腔”等文风问题交给独立的复核编辑会话，不用正则规则拦截——
        程序拦截文风既容易误杀，也会诱导模型用同义改写规避规则。
\end{itemize}

% ------------------------------------------------------------
\subsection{全文导读}
\label{sec:intro-roadmap}

表~\ref{tab:intro-roadmap} 给出各章内容与阅读建议。只关心系统设计思想的读者可以读第 \ref{sec:problem}、\ref{sec:principles}、
\ref{sec:arch} 节后跳到第 \ref{sec:guards}、\ref{sec:pipeline}、\ref{sec:eval} 节；关心金融方法的读者可以重点读第
\ref{sec:extract}～\ref{sec:reportcheck} 节；希望复现结果的读者可直接看附录~\ref{app:repro}。

\begin{table}[h]
\centering
\caption{全文导读。}
\label{tab:intro-roadmap}
\small
\begin{tabularx}{\linewidth}{@{}p{1.3cm}p{3.2cm}X@{}}
\toprule
章节 & 标题 & 内容 \\
\midrule
\S\ref{sec:related} & 背景与相关工作 & A 股年报结构、卖方研报体例与估值方法；财报问答、程序辅助推理、金融智能体、多智能体框架、自我改进、可追溯生成、LLM 评审与编码智能体 harness \\
\S\ref{sec:problem} & 问题形式化 & 统一记号；三项任务的输入、输出与正确性标准；可追溯性、数字闭合、可复现性；威胁模型与系统不变量；研报生成与评价的形式化 \\
\S\ref{sec:principles} & 设计原则 & 十一条设计规范的动机、不变量与执行机制 \\
\S\ref{sec:arch} & 系统架构 & 四层结构、模块划分、追踪与复现 \\
\S\ref{sec:extract} & 证据约束的双通道抽取 & 文档解析、关键页定位、规则通道、模型通道、逐字证据校验、合并与阻断 \\
\S\ref{sec:analysis} & 财报分析 & 真值表闭包、单季拆分、勾稽、异常信号、利润归因、一致预期与业绩预告、EPS 质量 \\
\S\ref{sec:valuation} & 估值建模 & 分部预测、三表联动、三表驱动 DCF、股权价值桥、敏感性、反向 DCF、可比公司、评级、工作簿 \\
\S\ref{sec:reportcheck} & 研报纠错 & 断言抽取、按展示精度的容差、错误类型诊断、内部一致性 \\
\S\ref{sec:guards} & 数字护栏与数字锁 & guard 的四项检查与出处分类、外部来源登记、numlock \\
\S\ref{sec:pipeline} & 多智能体研报流水线 & 隔离工作区、五个角色、11 道闸门、交接文件、refine 闭环 \\
\S\ref{sec:eval} & 评测体系 & 四层评分标准、盲评、计分、问题单路由、评审校准 \\
\S\ref{sec:exp} & 实验 & 抽取、纠错、估值一致性、研报质量、评审校准、成本 \\
\S\ref{sec:discussion}～\ref{sec:conclusion} & 讨论、局限、结论 & 设计取舍的回顾、已知问题与风险、后续工作 \\
附录 & & 复现说明（\ref{app:repro}）、评分标准全文（\ref{app:rubric}）、闸门定义（\ref{app:gates}）、数据格式（\ref{app:schemas}）、劣化算子（\ref{app:degrade}）、角色提示（\ref{app:prompts}）、第三方组件（\ref{app:thirdparty}） \\
\bottomrule
\end{tabularx}
\end{table}

关于记号：第 \ref{sec:problem} 节定义的记号在全文统一使用——年报 $\Doc$、页集合 $\Pages$、事实 $f=(m,k,v,u,\Src)$、
真值表 $\Truth$、研报断言集 $\Claims$、外部来源登记 $\Sources$、文稿 $\Narr$、本次运行的结论性产物 $\Art$、闸门 $\Gate_i$。
关于单位：财务事实内部一律以“元”存储；估值模块（三表、DCF、可比公司）以“百万元”计，每股以“元”计；
比率存为百分数值（14.49 表示 14.49\%）。文中引用的代码位置均相对于仓库根目录。
